\documentclass[a4paper,12pt]{article}

\usepackage[utf8]{inputenc}
\usepackage[brazil]{babel}
\usepackage{geometry}
\usepackage{graphicx}
\usepackage{caption}
\usepackage{subcaption}
\usepackage{titlesec}
\usepackage{float}
\usepackage[hidelinks]{hyperref}

\geometry{margin=2.5cm}
\graphicspath{{src/assets/images/}{./}}

\captionsetup{
    labelformat=empty,
    justification=centering,
    font=small,
    skip=10pt
}

\captionsetup[subfigure]{
    labelformat=empty,
    justification=centering,
    font=small
}

\titleformat{\section}
  {\large\bfseries}
  {}
  {0pt}
  {}

\title{\textbf{Game Design Document}}
\author{Juan Felipe Silva de Souza \\ João Pedro Silverio Gama}
\date{}

\begin{document}

\maketitle

%----------------------------------------

\section*{1. High Concept}

\noindent
\textit{Betrayal} é um roguelite 2D top-down de combate em arena, com progressão por sobrevivência, troca instantânea de armas elementais e escalonamento constante de dificuldade. O jogador controla um rei traído que precisa sobreviver ao cerco de soldados do usurpador, evoluir seus atributos ao longo da partida e derrotar o chefe final que surge no fim da run.

%----------------------------------------

\section*{2. Gameplay e Enredo}

\subsection*{Enredo}

O protagonista é um rei traído que teve seu trono usurpado. Cercado por tropas inimigas, ele atravessa uma arena medieval tomada por soldados, arqueiros e cavaleiros até enfrentar o rei usurpador, chefe final da partida.

\subsection*{Loop Principal}

\begin{itemize}
    \item O jogador inicia a partida no centro de uma arena maior que a tela, com câmera seguindo seu deslocamento.
    \item Inimigos surgem continuamente ao redor da área jogável e aumentam de pressão com o passar do tempo.
    \item O ataque é automático, baseado na espada equipada e em sua velocidade de ataque.
    \item Ao derrotar inimigos, o jogador recebe experiência, sobe de nível e escolhe uma entre três cartas de melhoria.
    \item A run termina quando o jogador morre ou derrota o chefe final, que aparece após aproximadamente 10 minutos de sobrevivência.
\end{itemize}

\pagebreak

\subsection*{Controles Atuais}

\begin{itemize}
    \item \textbf{WASD / Setas} -- movimentação do personagem.
    \item \textbf{SHIFT} -- esquiva com invulnerabilidade temporária.
    \item \textbf{SPACE} -- postura de guarda/estado defensivo.
    \item \textbf{1, 2 e 3} -- troca imediata entre as espadas de fogo, gelo e vento.
    \item \textbf{ESC} -- abre ou fecha o menu de pausa.
    \item \textbf{X} -- alterna a exibição de FPS para depuração.
\end{itemize}

\subsection*{Combate e Mobilidade}

O combate atual é focado em sobrevivência e posicionamento. Como os ataques acontecem automaticamente, a tomada de decisão do jogador se concentra em movimentação, esquiva, escolha da arma ideal para cada situação e gerenciamento de espaço contra grupos de inimigos e projéteis.

O personagem possui:

\begin{itemize}
    \item \textbf{Esquiva} -- deslocamento rápido com invulnerabilidade curta.
    \item \textbf{Regeneração passiva} -- cura gradual com base no atributo de recuperação.
    \item \textbf{Defesa} -- reduz parte do dano recebido.
\end{itemize}

\subsection*{Armas}

Diferente da versão inicial do GDD, o jogador começa a run com acesso às três espadas e pode alternar entre elas a qualquer momento:

\begin{itemize}
    \item \textbf{Espada de Fogo} -- causa bom dano base, aplica queimadura em área curta e é especialmente eficiente contra cavaleiros.
    \item \textbf{Espada de Gelo} -- desacelera inimigos, causa dano extra em combos consecutivos e pode congelar alvos após acertos sucessivos.
    \item \textbf{Espada de Vento} -- realiza golpes circulares em área, acerta mais de uma vez por ataque e aumenta a mobilidade do jogador.
\end{itemize}

\pagebreak

\subsection*{Progressão e Melhorias}

Cada subida de nível abre uma tela com três opções aleatórias dentre os atributos disponíveis. Cada atributo pode ser melhorado até cinco vezes durante a run.

\begin{itemize}
    \item \textbf{Vida Máxima} -- aumenta a vida total e concede uma pequena cura imediata.
    \item \textbf{Recuperação} -- melhora a regeneração passiva de vida.
    \item \textbf{Armadura} -- aumenta a redução de dano.
    \item \textbf{Força} -- aumenta o dano bruto dos ataques.
    \item \textbf{Velocidade de Movimento} -- melhora a mobilidade geral.
    \item \textbf{Velocidade de Ataque} -- reduz o intervalo entre golpes automáticos.
\end{itemize}

\subsection*{Mapa e Estrutura da Arena}

A arena atual não é uma tela fixa única. O mapa possui dimensões maiores que o viewport e inclui:

\begin{itemize}
    \item Câmera com seguimento do jogador;
    \item Obstáculos estáticos com colisão;
    \item Geração de ambiente com árvores;
    \item Uma vila com casas que interfere na navegação e no posicionamento durante o combate.
\end{itemize}

\pagebreak

%----------------------------------------

\section*{3. Inimigos e Escalonamento}

\subsection*{Tipos de Inimigo}

\begin{itemize}
    \item \textbf{Soldado} -- unidade corpo a corpo de pressão constante.
    \item \textbf{Arqueiro} -- inimigo de longo alcance que dispara flechas e exige deslocamento lateral.
    \item \textbf{Cavaleiro} -- inimigo resistente, com armadura elevada e comportamento de guarda.
    \item \textbf{Rei Usurpador} -- chefe final com grande vida, armadura própria e ciclo de ataques elementais inspirados nas três espadas do jogador.
\end{itemize}

\subsection*{Chefe Final}

O chefe aparece após cerca de 600 segundos de partida. Seu combate foi construído para funcionar como clímax da run:

\begin{itemize}
    \item Alterna ataques de vento, fogo e gelo;
    \item Usa golpes em área com tempo de preparação;
    \item Pode aplicar queimadura, lentidão e congelamento no jogador;
    \item Reduz o ritmo de surgimento de inimigos comuns enquanto está ativo;
    \item Ao ser derrotado, encerra a run com tela de vitória.
\end{itemize}

\subsection*{Escalonamento de Dificuldade}

O desafio não depende apenas de ondas fixas. A dificuldade cresce dinamicamente com:

\begin{itemize}
    \item Aumento progressivo da taxa de spawn ao longo do tempo;
    \item Crescimento de velocidade dos inimigos durante a partida;
    \item Aumento do dano inimigo conforme o tempo de sobrevivência e o nível do jogador.
\end{itemize}

\pagebreak

%----------------------------------------

\section*{4. Interface de Usuário}

\subsection*{Fluxo de Telas}

\begin{itemize}
    \item \textbf{Menu Principal} -- opções de jogar, abrir configurações e sair.
    \item \textbf{Cena de Jogo} -- combate principal com HUD e overlays de eventos.
    \item \textbf{Menu de Pausa} -- retomar, abrir opções ou voltar ao menu.
    \item \textbf{Configurações} -- controle de volume por slider.
    \item \textbf{Tela de Derrota} -- reinício rápido da run.
    \item \textbf{Tela de Vitória} -- reiniciar ou voltar ao menu principal.
\end{itemize}

\subsection*{HUD}

A interface em jogo foi simplificada e organizada em torno de leitura rápida de combate. Os elementos atualmente implementados são:

\begin{itemize}
    \item Barra de vida com valor atual e máximo;
    \item Barra de experiência com indicação de nível;
    \item Barra de armas com ícones clicáveis das três espadas;
    \item Contador de abates;
    \item Cronômetro da run;
    \item Exibição opcional de FPS.
\end{itemize}

\subsection*{Feedback Visual}

\begin{itemize}
    \item Números flutuantes de dano para golpes diretos e queimadura;
    \item Overlays de status para queimadura e congelamento;
    \item Cartas animadas de level up;
    \item Arcos e círculos de ataque desenhados na tela para indicar alcance e área dos golpes.
\end{itemize}

\pagebreak

%----------------------------------------

\section*{5. Áudio e Música}

O projeto já conta com uma trilha tema em loop durante a partida e com efeitos sonoros separados para cada arma elemental. A direção de áudio reforça o ritmo da run e ajuda o jogador a identificar a espada equipada no momento.

\begin{itemize}
    \item Música tema tocada durante a run;
    \item Efeito sonoro exclusivo para cada golpe de fogo, gelo e vento.
\end{itemize}

%----------------------------------------

\section*{6. Direção de Arte e Referências}

O estilo visual de \textit{Betrayal} segue a linha de pixel art 2D em perspectiva top-down, com temática medieval-fantástica. A apresentação atual combina sprites de personagens, armas e obstáculos ambientais com UI de forte contraste para priorizar leitura durante combates intensos.

As principais referências de gênero continuam próximas de jogos como \textit{Vampire Survivors}, \textit{The Binding of Isaac} e \textit{Enter the Gungeon}, mas com foco próprio em espadas elementais e sobrevivência em arena.

%----------------------------------------

\section*{7. Créditos de Assets}

\subsection*{Sprites}

\begin{itemize}
    \item \textbf{Knight} -- PixiVan (\href{https://pixivan.itch.io/}{pixivan.itch.io})
    \item \textbf{Player} -- SventHole (\href{https://sventhole.itch.io/}{sventhole.itch.io})
    \item \textbf{Armor} -- Polygonix (\href{https://polygonix.itch.io/}{polygonix.itch.io})
    \item \textbf{King} -- LuizMelo (\href{https://luizmelo.itch.io/}{luizmelo.itch.io})
    \item \textbf{Archer and Soldier} -- Zerie (\href{https://zerie.itch.io}{zerie.itch.io})
\end{itemize}

\subsection*{Sound Effects}

\begin{itemize}
    \item \textbf{Fire Slash} -- Fronbondi\_Skegs (\href{https://pixabay.com/users/fronbondi_skegs-23154649/}{pixabay.com})
    \item \textbf{Ice Slash} -- Abdullah Pullin (\href{https://samplefocus.com/users/abdullah-pullin}{samplefocus.com})
    \item \textbf{Wind Slash} -- Jesse Nicasio (\href{https://samplefocus.com/users/jesse-nicasio}{samplefocus.com})
\end{itemize}

%----------------------------------------

\section*{Imagens do Projeto}

\begin{figure}[H]
    \centering
    
\includegraphics[width=0.55\textwidth]{logo.png}
    \caption{\textit{Logo atual do jogo Betrayal}}
\end{figure}

\begin{figure}[H]
    \centering
    \begin{subfigure}[t]{0.47\textwidth}
        \centering
        \includegraphics[width=\textwidth]{player_spritesheet.png}
        \caption{\textit{Spritesheet do protagonista controlado pelo jogador}}
    \end{subfigure}
    \hfill
    \begin{subfigure}[t]{0.47\textwidth}
        \centering
        \includegraphics[width=\textwidth]{soldier_spritesheet.png}
        \caption{\textit{Spritesheet base utilizada para soldados e arqueiros}}
    \end{subfigure}

    \vspace{0.5cm}

    \begin{subfigure}[t]{0.47\textwidth}
        \centering
        \includegraphics[width=\textwidth]{knight_spritesheet.png}
        \caption{\textit{Spritesheet do cavaleiro blindado}}
    \end{subfigure}
    \hfill
    \begin{subfigure}[t]{0.47\textwidth}
        \centering
        \includegraphics[width=\textwidth]{boss_spritesheet.png}
        \caption{\textit{Spritesheet do chefe final}}
    \end{subfigure}
\end{figure}

\begin{figure}[H]
    \centering
    \begin{subfigure}[t]{0.28\textwidth}
        \centering
        
\includegraphics[width=0.95\textwidth]{fire_sword.png}
        \caption{\textit{Espada de Fogo}}
    \end{subfigure}
    \hfill
    \begin{subfigure}[t]{0.28\textwidth}
        \centering
        
\includegraphics[width=0.95\textwidth]{ice_sword.png}
        \caption{\textit{Espada de Gelo}}
    \end{subfigure}
    \hfill
    \begin{subfigure}[t]{0.28\textwidth}
        \centering
        \includegraphics[width=0.95\textwidth]{wind_sword.png}
        \caption{\textit{Espada de Vento}}
    \end{subfigure}
\end{figure}

\begin{figure}[H]
    \centering
    \begin{subfigure}[t]{0.28\textwidth}
        \centering
        
\includegraphics[width=0.95\textwidth]{upgrades/health_up.png}
        \caption{\textit{Vida Máxima}}
    \end{subfigure}
    \hfill
    \begin{subfigure}[t]{0.28\textwidth}
        \centering
        
\includegraphics[width=0.95\textwidth]{upgrades/regen.png}
        \caption{\textit{Recuperação}}
    \end{subfigure}
    \hfill
    \begin{subfigure}[t]{0.28\textwidth}
        \centering
        
\includegraphics[width=0.95\textwidth]{upgrades/armor.png}
        \caption{\textit{Armadura}}
    \end{subfigure}

    \vspace{0.5cm}

    \begin{subfigure}[t]{0.28\textwidth}
        \centering
        
\includegraphics[width=0.95\textwidth]{upgrades/strength.png}
        \caption{\textit{Força}}
    \end{subfigure}
    \hfill
    \begin{subfigure}[t]{0.28\textwidth}
        \centering
        
\includegraphics[width=0.95\textwidth]{upgrades/move_speed.png}
        \caption{\textit{Velocidade de Movimento}}
    \end{subfigure}
    \hfill
    \begin{subfigure}[t]{0.28\textwidth}
        \centering
        
\includegraphics[width=0.95\textwidth]{upgrades/attack_speed.png}
        \caption{\textit{Velocidade de Ataque}}
    \end{subfigure}
\end{figure}

\begin{figure}[H]
    \centering
    \includegraphics[width=0.82\textwidth]{buildings.png}
    \caption{\textit{Spritesheet de construções utilizadas na composição da arena}}
\end{figure}

\end{document}
